\section{Methods}

\subsection{Task and label space}

We study chest X-ray (CXR) vision-language models on structured multi-label finding prediction. The primary label space contains eight findings:

\begin{verbatim}
atelectasis
cardiomegaly
consolidation
edema
pleural_effusion
pneumonia
pneumothorax
no_finding
\end{verbatim}

Each model outputs a structured JSON object containing:

\begin{itemize}
\item \texttt{findings[label] $\in$ \{positive, negative, uncertain, not\_mentioned\}};
\item \texttt{probability\_positive[label] $\in$ [0,1]};
\item an optional short rationale.
\end{itemize}

All models share one prompt and one parser. Missing labels are mapped to \texttt{not\_mentioned}; invalid statuses are mapped to \texttt{uncertain}; invalid probabilities are set to null. Parse success and parse failure are recorded separately.

\subsection{Models}

We evaluate openly available vision-language models:

\begin{table}[t]
\centering
\small
\begin{tabular}{ll}
\hline
Model & Role \\
MedGemma-4B-IT & medical generalist primary model \\
Qwen2.5-VL-7B-Instruct & open-weight general baseline \\
CheXagent-8B & CXR-specific pilot/sensitivity model \\
\hline
\end{tabular}
\end{table}


Models are loaded locally with 4-bit quantization and a single-GPU placement when available. MedGemma and Qwen are used for the main Stage 2 experiments; CheXagent is used only for a small pilot because it is substantially slower under CPU offload.

\subsection{Data}

\subsubsection{NIH ChestX-ray14 pilot}

We use a stratified pilot of 1,000 studies from NIH ChestX-ray14. The manifest contains patient-level splits:

\begin{verbatim}
internal_train
source_val
internal_test
\end{verbatim}

\texttt{source\_val} is used only for calibration fitting and threshold selection. \texttt{internal\_test} is used only for evaluation. Image-level ground-truth labels are derived from the official NIH labels via a fixed mapping into the primary label space.

\subsubsection{CheXpert-plus-RRG external report set}

CheXpert-plus-RRG contains:

\begin{verbatim}
findings_section:  62 samples
impression_section: 202 samples
\end{verbatim}

Each sample provides a frontal CXR image and reference report text. It is used as an external report-grounded evaluation set for hallucination, omission, and image reliance, not as a large-scale classification benchmark.

\subsection{Structured output and parsing}

The structured prompt asks the model to return one JSON object with findings and probabilities. The parser:

\begin{enumerate}
\item extracts the first JSON object;
\item tolerates Markdown code fences;
\item normalises misspelled label keys with fuzzy matching;
\item maps missing labels to \texttt{not\_mentioned};
\item records parse failure explicitly.
\end{enumerate}

\subsection{Distribution shift benchmark}

We evaluate six image-quality shifts at three severity levels:

\begin{table}[t]
\centering
\small
\begin{tabular}{llll}
\hline
Shift & Mild & Moderate & Severe \\
Blur & $\sigma$=0.5 & $\sigma$=1.0 & $\sigma$=2.0 \\
Noise & $\sigma$=0.02 & $\sigma$=0.05 & $\sigma$=0.10 \\
JPEG & q=70 & q=40 & q=20 \\
Brightness & $\times$0.8 & $\times$1.2 & $\times$1.5 \\
Contrast & $\times$0.7 & $\times$1.3 & $\times$1.6 \\
Resolution & $\times$0.5 & $\times$0.25 & $\times$0.125 \\
\hline
\end{tabular}
\end{table}


Each image is evaluated under clean, mild, moderate, and severe conditions. For every image--condition--severity combination we store the raw output, parsed output, parse success, probability estimates, elapsed time, and ground truth.

\subsection{Image-reliance evaluation}

For image reliance we compare the original image against:

\begin{itemize}
\item blank gray image;
\item strong blur;
\item shuffled image from the same manifest.
\end{itemize}

For each condition and label we measure:

\begin{itemize}
\item status flip count relative to original;
\item image sensitivity (fraction of images with at least one finding flip);
\item mean absolute change in \texttt{probability\_positive}.
\end{itemize}

\subsection{Calibration}

For each finding we fit temperature scaling on \texttt{source\_val} and apply it to \texttt{internal\_test} without refitting. We report:

\begin{itemize}
\item AUROC;
\item expected calibration error (ECE);
\item Brier score.
\end{itemize}

Temperature scaling changes calibration but not ranking, so AUROC is invariant by construction.

\subsection{Selective prediction}

For each finding prediction we derive a confidence score:

\begin{verbatim}
confidence = P(positive)        if predicted positive
confidence = 1 - P(positive)    if predicted negative
confidence = 0.5                if uncertain
risk = 1 - confidence
\end{verbatim}

We compute risk--coverage curves by accepting the lowest-risk cases first, and report AURC, coverage at 5\% and 10\% risk, and selective accuracy at 80\% coverage.

\subsection{CheXpert-RRG report-level evaluation}

RRG reference text is converted to finding-level evidence using a rule-based keyword + negation-window method (\texttt{rule\_based\_v0}). For each primary finding we define:

\begin{itemize}
\item disease hallucination: model positive while reference evidence is not present;
\item negation failure: model positive while reference explicitly negates the finding;
\item omission: model negative while reference evidence is present.
\end{itemize}

This evidence extractor is a pilot implementation; final analysis should use a validated report labeller or human adjudication.

\subsection{Statistics}

We use image-level cluster bootstrap to estimate 95\% confidence intervals. For paired clean-vs-shift comparisons we resample images with replacement and compute paired bootstrap differences in macro AUROC, ECE, and Brier score. The current Stage 2 analysis uses 200 bootstrap resamples for pilot speed; the final paper should use at least 1,000.
